% ==================== BAB IV ====================
\section{Metodologi dan Cara Kerja Sistem}
\label{bab:4}

\subsection{Diagram Blok Sistem Kendali}
Sistem kendali gerak presisi tinggi yang dirancang pada penelitian ini mengadopsi topologi lingkar tertutup (\textit{closed-loop control system}). Sistem ini memastikan aktuator linier bergerak dengan presisi tinggi melalui mekanisme umpan balik (\textit{feedback}) yang secara kontinu mengoreksi deviasi posisi. Secara garis besar, arsitektur perangkat keras dibagi menjadi dua domain utama, yaitu \textit{Digital Controller} dan \textit{Servo Power System (Plant)}.

\begin{figure}[htbp]
    \centering
    % Pastikan nama file gambarnya sesuai dengan yang di folder lu
    \includegraphics[width=0.95\textwidth]{Diagram Blok Sistem.png}
    \caption{Diagram Blok Sistem Kendali Gerak \textit{Closed-Loop}}
    \label{fig:diagram_blok}
\end{figure}

Sesuai dengan alur pada Gambar \ref{fig:diagram_blok}, prinsip kerja sistem secara fungsional diuraikan sebagai berikut:
\begin{enumerate}
    \item \textbf{Input Setpoint Source (PC):} Berperan sebagai antarmuka pengguna (\textit{Human-Machine Interface}) untuk memasukkan parameter target posisi dan kecepatan yang diinginkan.
    \item \textbf{Summing Junction (Comparator):} Titik komparasi matematis di mana nilai target posisi (referensi) dikurangi dengan nilai posisi aktual (\textit{feedback signal}) yang dibaca dari \textit{encoder} motor. Selisih dari kedua nilai ini menghasilkan \textbf{Sinyal Error} yang kemudian dikirim ke unit pengontrol.
    \item \textbf{Digital Controller (ESP32-based Control System):} Otak utama dari sistem kendali yang memanfaatkan arsitektur \textit{dual-core}:
    \begin{itemize}
        \item \textbf{ESP32 (Core 0):} Bertugas memproses sinyal error dan melakukan kalkulasi matriks algoritma \textit{trajectory planning} (\textit{Rectangular, Trapezoidal}, atau \textit{S-Curve}). Hasil komputasi berupa parameter jeda waktu (\textit{delay}) dikirimkan ke inti kedua.
        \item \textbf{ESP32 (Core 1):} Berfungsi murni sebagai generator pulsa (\textit{pulse generator}) yang mengeksekusi instruksi dari Core 0 untuk membangkitkan sinyal kendali frekuensi tinggi.
    \end{itemize}
    \item \textbf{Sirkuit Antarmuka Transmisi:} Sinyal TTL dari ESP32 diamankan melalui \textbf{Isolasi Galvanik (PC817)} untuk mencegah kerusakan akibat arus balik. Sinyal tersebut kemudian diteruskan ke \textbf{Differential Line Driver (AM26LS31)} untuk dikonversi menjadi \textit{Differential Actuating Signal} berstandar RS-422 yang kebal terhadap derau elektromagnetik lingkungan industri.
    \item \textbf{Servo Power System (Plant):} \textbf{AC Servo Amplifier MR-J2S-10A} menerima sinyal diferensial tersebut dan memodulasikan daya AC untuk menggerakkan \textbf{Motor AC Servo (HC-KFS053)} pada \textbf{Aktuator Linier ED-4270m}. 
    \item \textbf{Instrumen Evaluasi:} Respons transien dan getaran sisa (\textit{settling time}) selama motor beroperasi dievaluasi secara \textit{real-time} dengan menghubungkan \textbf{Osiloskop} ke terminal sinyal analog (MO1) pada \textit{servo amplifier}.
\end{enumerate}

\subsection{Diagram Alir Sistem Perangkat Lunak (\textit{Flowchart})}
Kompleksitas eksekusi algoritma gerak pada mikrokontroler \textit{dual-core} dipetakan dalam sebuah diagram alir sistem yang terbagi menjadi dua bagian berkesinambungan.

\begin{figure}[htbp]
    \centering
    \begin{minipage}{0.45\textwidth}
        \centering
        \includegraphics[width=\linewidth]{Diagram Alir Sistem A.png}
    \end{minipage}\hfill
    \begin{minipage}{0.35\textwidth}
        \centering
        \includegraphics[width=\linewidth]{Diagram Alir Sistem B.png}
    \end{minipage}
    \caption{Diagram Alir Sistem Perangkat Lunak Keseluruhan}
    \label{fig:flowchart_sistem}
\end{figure}

Berdasarkan Gambar \ref{fig:flowchart_sistem}, alur kerja peranti lunak (\textit{firmware}) beroperasi dengan tahapan sebagai berikut:
\begin{enumerate}
    \item \textbf{Inisialisasi dan Komunikasi:} Program dimulai dengan inisialisasi pin I/O, pengaturan \textit{Hardware Timer}, dan pengaktifan sistem RTOS. Sistem kemudian bersiap menerima data parameter operasional dari antarmuka PC (HMI).
    \item \textbf{Kalkulasi Trajektori (Core 0):} Setelah data diterima, \textit{Core 0} memproses matriks \textit{trajectory planning} sesuai dengan profil gerak yang dipilih pengguna (\textit{Rectangular, Trapezoidal}, atau \textit{S-Curve}). Keluaran dari proses ini adalah parameter jeda waktu transmisi pulsa.
    \item \textbf{Eksekusi Interupsi (Core 1):} Memasuki konektor logika transisi, \textit{Core 1} mengambil alih tugas aktuasi dengan mengeksekusi \textit{Hardware Timer Interrupt}. Inti ini membangkitkan dan mengirimkan sinyal pulsa menuju sirkuit transmisi diferensial secara presisi.
    \item \textbf{Pergerakan dan Pemantauan Target:} Aktuator motor bergerak menunju target posisi yang diinstruksikan (Konektor A). Sistem secara terus-menerus memantau apakah target posisi akhir telah tercapai.
    \item \textbf{Looping atau Terminasi:} Jika target posisi belum tercapai, logika program akan berbalik (Konektor B) untuk memerintahkan \textit{Core 1} agar mengulangi siklus penembakan pulsa. Jika target telah tercapai, mikrokontroler akan menghentikan pulsa dan motor berhenti seketika.
    \item \textbf{Evaluasi Kinerja:} Pada fase akhir, pengujian kualitas gerak dilakukan dengan mengevaluasi durasi \textit{settling time} dan observasi efek \textit{liquid slosh} menggunakan instrumen osiloskop sebelum program dinyatakan selesai.
\end{enumerate}